\par\medskip
\SourceLeaf{078}\pdfbookmark[2]{59.}{article-059}
\begin{center}59.\end{center}
The rules for determining distances from planes passing through the Sun can clearly also be applied to the distances of the Earth. Here, however, usually only the simplest cases arise. Let $R$ be the distance of the Earth from the Sun, $L$ the heliocentric longitude of the Earth (which differs by $180^0$ from the geocentric longitude of the Sun), and $X$, $Y$, $Z$ the distances of the Earth from three planes intersecting at right angles at the Sun. We now consider the following cases.

I. If the plane for the distances $Z$ is the ecliptic itself, and the longitudes of the poles of the other planes, from which the distances are $X$, $Y$, are $N$ and $N+90^0$ respectively, then
\[
X=R\cos(L-N),\qquad Y=R\sin(L-N),\qquad Z=0.
\]

II. If the plane for the distances $Z$ is parallel to the equator, and the right ascensions of the poles of the other planes, from which the distances are $X$, $Y$, are $0$ and $90^0$ respectively, we have, with the obliquity of the ecliptic denoted by $\epsilon$,
\[
X=R\cos L,\qquad Y=R\cos\epsilon\sin L,\qquad Z=R\sin\epsilon\sin L.
\]
The distinguished \textsc{de Zach} and \textsc{de Lambre}, editors of the most recent solar tables, have begun to take the latitude of the Sun into account. This quantity, produced by the perturbations of the other planets and the Moon, can scarcely reach one arcsecond. We denote the heliocentric latitude of the Earth, which always agrees in magnitude with the Sun's latitude and has the opposite sign, by $B$. We then have
\[
\begin{array}{@{}l|l@{}}
\text{in case I.}&\text{in case II.}\\[.4em]
X=R\cos B\cos(L-N)&X=R\cos B\cos L\\[.4em]
Y=R\cos B\sin(L-N)&Y=R\cos B\cos\epsilon\sin L-R\sin B\sin\epsilon\\[.4em]
Z=R\sin B&Z=R\cos B\sin\epsilon\sin L+R\sin B\cos\epsilon.
\end{array}
\]
Here $\cos B$ can always safely be replaced by $1$, and the angle $B$ expressed in parts of the radius can be used for $\sin B$.

The coordinates thus obtained refer to the \textit{center} of the Earth. Let $\xi$, $\eta$, $\zeta$ be the distances of any point on the Earth's \textit{surface} from three planes passing through the Earth's center and parallel to those passing through the Sun. The distances of that point from the \SourceLeaf{079}planes passing through the Sun will clearly be $X+\xi$, $Y+\eta$, $Z+\zeta$. The values of the coordinates $\xi$, $\eta$, $\zeta$ are readily determined in both cases as follows. Let $\rho$ be the radius of the terrestrial globe (or the sine of the mean horizontal parallax of the Sun). Let $\lambda$ be the longitude of the point on the celestial sphere onto which the line drawn from the Earth's center to the point on its surface is projected, $\beta$ its latitude, $\alpha$ its right ascension, and $\delta$ its declination. Then
\[
\begin{array}{@{}l|l@{}}
\text{in case I.}&\text{in case II.}\\[.4em]
\xi=\rho\cos\beta\cos(\lambda-N)&\xi=\rho\cos\delta\cos\alpha\\[.4em]
\eta=\rho\cos\beta\sin(\lambda-N)&\eta=\rho\cos\delta\sin\alpha\\[.4em]
\zeta=\rho\sin\beta&\zeta=\rho\sin\delta.
\end{array}
\]
That point on the celestial sphere clearly coincides with the zenith of the site on the surface if the Earth is treated as a sphere. Its right ascension will therefore agree with the right ascension of the midheaven, or the sidereal time converted into degrees, and its declination with the elevation of the pole. If taking the spheroidal shape of the Earth into account were worth the effort, one would use for $\delta$ the \textit{corrected} elevation of the pole and for $\rho$ the true distance of the site from the Earth's center. Both are obtained by known rules. From $\alpha$ and $\delta$ the longitude and latitude $\lambda$ and $\beta$ are obtained by the known rules that will also be given below. It is also evident that $\lambda$ agrees with the longitude of the \textit{nonagesimal point}, and $90^0-\beta$ with its altitude.
